\documentclass[11pt,a4paper]{article}
\usepackage[utf8]{inputenc}
\usepackage[margin=1in]{geometry}
\usepackage{amsmath,amssymb,amsthm}
\usepackage{listings}
\usepackage{xcolor}
\usepackage{hyperref}

\newtheorem{theorem}{Theorem}
\newtheorem{lemma}[theorem]{Lemma}

\lstset{
  language=C++,
  basicstyle=\ttfamily\small,
  keywordstyle=\color{blue}\bfseries,
  commentstyle=\color{gray},
  stringstyle=\color{red},
  numbers=left,
  numberstyle=\tiny\color{gray},
  breaklines=true,
  frame=single,
  tabsize=4,
  showstringspaces=false
}

\title{IOI 2016 -- Messy}
\author{}
\date{}

\begin{document}
\maketitle

\section{Problem Summary}

Given $n$ (a power of 2), determine an unknown permutation $P$ of $\{0, 1, \ldots, n-1\}$ in two phases:
\begin{itemize}
  \item \textbf{Phase 1 (Add)}: Insert binary strings of length $n$ into a set $S$.
  \item \textbf{Shuffle}: $P$ is applied to every string in $S$ (bit $i$ moves to position $P(i)$).
  \item \textbf{Phase 2 (Check)}: Query whether a binary string belongs to the shuffled set.
\end{itemize}
Both phases allow at most $O(n \log n)$ operations.

\section{Solution}

Use \textbf{divide and conquer on bit positions}.

\subsection{Phase 1: Adding Strings}

For D\&C on range $[l, r)$ with a ``context'' set $C$ (bit positions outside $[l, r)$ that are always set to~1):
\begin{enumerate}
  \item Let $\text{mid} = (l + r) / 2$.
  \item For each $i \in [l, \text{mid})$: add a string with bit $i$ set to~1 and all bits in $C$ set to~1 (all others~0).
  \item Recurse on $[l, \text{mid})$ with $C' = C \cup [\text{mid}, r)$, and on $[\text{mid}, r)$ with $C' = C \cup [l, \text{mid})$.
\end{enumerate}

Each level adds $n/2$ strings across all recursive calls, totalling $O(n \log n)$ strings over $\log n$ levels.

\subsection{Phase 2: Checking Membership}

After the shuffle, for D\&C on range $[l, r)$ with a known set of candidate output positions and shuffled context bits:
\begin{enumerate}
  \item For each candidate position $p$, check the string with bit $p$ and context bits set. If present, $p$ maps to the left half $[l, \text{mid})$; otherwise, to the right half.
  \item Recurse, using the classified positions as context for sub-problems.
\end{enumerate}

Each level performs $n$ checks, totalling $O(n \log n)$.

\section{C++ Implementation}

\begin{lstlisting}
#include <bits/stdc++.h>
using namespace std;

// Grader: void add_element(string), bool check_element(string),
//         void compile_set(), void answer(int[])

int n;
int P[1024];

void addStrings(int l, int r, vector<int> &ctx) {
    if (r - l <= 1) return;
    int mid = (l + r) / 2;
    for (int i = l; i < mid; i++) {
        string s(n, '0');
        s[i] = '1';
        for (int c : ctx) s[c] = '1';
        add_element(s);
    }
    vector<int> leftCtx = ctx;
    for (int i = mid; i < r; i++) leftCtx.push_back(i);
    addStrings(l, mid, leftCtx);

    vector<int> rightCtx = ctx;
    for (int i = l; i < mid; i++) rightCtx.push_back(i);
    addStrings(mid, r, rightCtx);
}

void findPerm(int l, int r, vector<int> &positions, vector<int> &ctx) {
    if (r - l == 1) { P[l] = positions[0]; return; }
    int mid = (l + r) / 2;
    vector<int> leftPos, rightPos;
    for (int p : positions) {
        string s(n, '0');
        s[p] = '1';
        for (int c : ctx) s[c] = '1';
        if (check_element(s)) leftPos.push_back(p);
        else rightPos.push_back(p);
    }
    vector<int> leftCtx = ctx;
    for (int p : rightPos) leftCtx.push_back(p);
    findPerm(l, mid, leftPos, leftCtx);

    vector<int> rightCtx = ctx;
    for (int p : leftPos) rightCtx.push_back(p);
    findPerm(mid, r, rightPos, rightCtx);
}

void restore_permutation(int N, int w, int r) {
    n = N;
    vector<int> emptyCtx;
    addStrings(0, n, emptyCtx);
    compile_set();
    vector<int> allPos(n);
    iota(allPos.begin(), allPos.end(), 0);
    vector<int> emptyCtx2;
    findPerm(0, n, allPos, emptyCtx2);
    answer(P);
}
\end{lstlisting}

\section{Complexity Analysis}

\begin{itemize}
  \item \textbf{Add operations}: $O(n \log n)$.
  \item \textbf{Check operations}: $O(n \log n)$.
  \item \textbf{Time}: $O(n^2 \log n)$ total work (each string operation touches $O(n)$ bits).
  \item \textbf{Space}: $O(n \log n)$ for the set of strings.
\end{itemize}

\end{document}
